% !TEX program = lualatex
% =====================================================================
%  JOY2Mulka 利用マニュアル
%  ビルド:  lualatex JOY2Mulka_Manual.tex   （2回実行して目次を確定）
% =====================================================================
\documentclass[a4paper,11pt]{ltjsarticle}

\usepackage{geometry}
\usepackage{graphicx}
\usepackage{xcolor}
\usepackage{colortbl}
\usepackage{tcolorbox}
\usepackage{fancyhdr}
\usepackage{enumitem}
\usepackage{booktabs}
\usepackage{array}
\usepackage{longtable}
\usepackage{needspace}
\usepackage[hidelinks]{hyperref}

\tcbuselibrary{skins,breakable}

\geometry{top=22mm,bottom=22mm,left=20mm,right=20mm,headsep=6mm,footskip=12mm}

\definecolor{brand}{RGB}{37,99,235}      % アプリのプライマリ（青）
\definecolor{brandink}{RGB}{23,42,79}
\definecolor{soft}{RGB}{240,245,252}
\definecolor{warn}{RGB}{180,120,0}
\definecolor{warnbg}{RGB}{254,249,231}
\definecolor{okgreen}{RGB}{22,124,80}
\definecolor{okbg}{RGB}{236,250,242}
\definecolor{grayrule}{gray}{0.75}
\definecolor{band}{gray}{0.92}

\pagestyle{fancy}
\fancyhf{}
\fancyhead[L]{\small\color{brandink}JOY2Mulka 利用マニュアル}
\fancyhead[R]{\small\color{brandink}\leftmark}
\fancyfoot[C]{\small\thepage}
\renewcommand{\headrulewidth}{0.4pt}
\renewcommand{\headrule}{\hbox to\headwidth{\color{grayrule}\leaders\hrule height \headrulewidth\hfill}}
\renewcommand{\sectionmark}[1]{\markboth{#1}{}}

\setlength{\parindent}{0pt}
\setlength{\parskip}{5pt}

% --- 見出しの体裁 -----------------------------------------------------
\makeatletter
\renewcommand{\section}{\@startsection{section}{1}{\z@}%
  {2.2\baselineskip \@plus .3ex}{0.9\baselineskip}%
  {\normalfont\Large\bfseries\color{brandink}}}
\renewcommand{\subsection}{\@startsection{subsection}{2}{\z@}%
  {1.5\baselineskip \@plus .2ex}{0.6\baselineskip}%
  {\normalfont\large\bfseries\color{brandink}}}
\renewcommand{\subsubsection}{\@startsection{subsubsection}{3}{\z@}%
  {1.1\baselineskip \@plus .2ex}{0.4\baselineskip}%
  {\normalfont\normalsize\bfseries\color{brandink}}}
\makeatother

% --- 便利マクロ -------------------------------------------------------
% 画面上のボタン・項目名
\newcommand{\ui}[1]{\textcolor{brand}{\textbf{［#1］}}}
% ファイル名・コマンド
\newcommand{\fn}[1]{\texttt{#1}}

% 図（幅指定つき、枠線あり）
\newcommand{\shot}[3][0.96]{%
  \begin{center}
  \fcolorbox{grayrule}{white}{\includegraphics[width=#1\textwidth]{img/#2}}\\[3pt]
  {\small\color{brandink}\textbf{図}\quad #3}
  \end{center}%
}

% 補足ボックス
\newtcolorbox{tipbox}[1][ヒント]{
  breakable, enhanced, colback=soft, colframe=brand,
  boxrule=0pt, leftrule=3pt, arc=1mm, left=8pt, right=8pt, top=6pt, bottom=6pt,
  fonttitle=\bfseries, title=#1, coltitle=brandink,
  attach boxed title to top left={xshift=8pt, yshift=-3pt},
  boxed title style={colback=soft, boxrule=0pt, sharp corners}
}
\newtcolorbox{warnbox}[1][注意]{
  breakable, enhanced, colback=warnbg, colframe=warn,
  boxrule=0pt, leftrule=3pt, arc=1mm, left=8pt, right=8pt, top=6pt, bottom=6pt,
  fonttitle=\bfseries, title=#1, coltitle=warn,
  attach boxed title to top left={xshift=8pt, yshift=-3pt},
  boxed title style={colback=warnbg, boxrule=0pt, sharp corners}
}
\newtcolorbox{stepbox}[1][操作手順]{
  breakable, enhanced, colback=okbg, colframe=okgreen,
  boxrule=0pt, leftrule=3pt, arc=1mm, left=8pt, right=8pt, top=6pt, bottom=6pt,
  fonttitle=\bfseries, title=#1, coltitle=okgreen,
  attach boxed title to top left={xshift=8pt, yshift=-3pt},
  boxed title style={colback=okbg, boxrule=0pt, sharp corners}
}

\setlist[enumerate]{itemsep=2pt, topsep=3pt, leftmargin=1.6em}
\setlist[itemize]{itemsep=2pt, topsep=3pt, leftmargin=1.4em}

\renewcommand{\arraystretch}{1.25}
\setlength{\tabcolsep}{5pt}

\begin{document}
\gtfamily\sffamily

% =====================================================================
\begin{titlepage}
\centering
\vspace*{40mm}
{\Huge\bfseries\color{brandink} JOY2Mulka}\\[6mm]
{\LARGE 利用マニュアル}\\[4mm]
{\large オリエンテーリング大会のスタートリスト生成ツール}\\[18mm]

\begin{tcolorbox}[width=0.82\textwidth, colback=soft, colframe=brand,
  boxrule=0pt, leftrule=4pt, arc=1mm, left=14pt, right=14pt, top=12pt, bottom=12pt]
このマニュアルは、\textbf{初めて JOY2Mulka を使う人が、読むだけで一通りの操作ができるようになる}ことを
目的に書かれています。画面のスクリーンショットを見ながら、上から順に読み進めてください。\\[4pt]
急いでいる場合は、\textbf{第\ref{sec:quick}章「最短手順」}だけを読めば、
エントリーリストからスタートリスト一式を生成できます。
\end{tcolorbox}

\vfill
{\small 対象バージョン: JOY2Mulka Web / デスクトップ版\quad ｜\quad 本書の作成日: 2026年9月9日}
\end{titlepage}

% =====================================================================
\tableofcontents
\clearpage

% =====================================================================
\section{JOY2Mulka とは}

JOY2Mulka は、オリエンテーリング大会の\textbf{エントリーリスト}（申込データ）を読み込んで、
運営に必要な\textbf{スタートリスト}を一式まとめて作るツールです。

手作業でやると半日かかる「クラスごとの並び替え」「同じクラブが連続しないような調整」
「スタート時刻の割り当て」「掲示用の清書」を、数分で終わらせることができます。

\subsection{何ができるのか（機能一覧）}

\begin{center}
\begin{tabular}{@{}>{\bfseries}p{42mm}p{105mm}@{}}
\toprule
\rowcolor{band}\textbf{機能} & \textbf{説明} \\
\midrule
エントリー読み込み & JOY（Japan-O-entrY）形式の CSV / XLSX を自動判別。
それ以外の形式でも、列の対応を画面上で指定すれば読み込める。 \\
クラス分割 & 人数の多いクラス（例: M21E 149名）を M21E1 / M21E2 のように複数コースへ均等分割。 \\
ランキング順分割 & JOA ランキングを使い、強い選手が偏らないようにグループ分け。 \\
複数スタート・レーン & スタートエリアとレーンをいくつでも作り、レーンごとに開始時刻・
スタート間隔・コース間の空きを設定。 \\
同一クラブ連続回避 & 同じ所属の選手が連続してスタートしないように並び替え。 \\
同時刻クラブ重複検出 & 生成後に「同じ時刻に同じ所属」「同じレーンの近い時刻に同じ所属」を検出して警告。 \\
人物位置制約 & 特定の人を「前半（先頭20\%）」「後半（最後20\%）」に寄せる。 \\
JOA 競技者番号取得 & 番号が空欄の選手について、JOA 登録者リストから自動補完（デスクトップ版のみ）。 \\
練習会モード & スタート時刻を割り当てず、入力順のままクラスごとに出力。 \\
ゼッケン番号の省略 & スタートナンバー列を出力から丸ごと外す。 \\
スタートリスト修正 & 生成済みのリストを読み込み、出走順・時刻の入れ替えなどを画面上で修正。 \\
形式変換 & CSV / TeX のスタートリストを、他の全形式へ一括変換。 \\
7種類の出力 & Mulka 用 CSV、役員用 CSV、公開用 / 役員用の LaTeX と Word、クラス別人数集計。 \\
\bottomrule
\end{tabular}
\end{center}

\subsection{生成されるファイル}
\label{sec:files}

最後の画面で、以下の 7 ファイルを個別に、あるいは ZIP でまとめてダウンロードできます。

\begin{center}
\begin{tabular}{@{}>{\ttfamily}p{45mm}p{50mm}p{45mm}@{}}
\toprule
\rowcolor{band}\textbf{\rmfamily ファイル名} & \textbf{用途} & \textbf{渡す相手・使い方} \\
\midrule
Startlist.csv & Mulka へのインポート用 & 計時担当。Mulka でそのまま読み込む \\
Role\_Startlist.csv & 役員用（チェックイン欄つき） & スタート地区の係員。印刷して受付に使う \\
Public\_Startlist.tex & 公開用スタートリスト & LuaLaTeX で PDF 化して掲示・Web 公開 \\
Role\_Startlist.tex & 役員用（ふりがな付き） & 呼び出し担当。読み仮名つきで名前を間違えない \\
Public\_Startlist.docx & 公開用（Word 形式） & LaTeX を使わない場合はこちら \\
Role\_Startlist.docx & 役員用（Word 形式） & 同上 \\
Class\_Summary.csv & クラス別人数の集計 & 地図の必要枚数、賞品数の確認 \\
\bottomrule
\end{tabular}
\end{center}

\subsection{3つのモード}

起動すると、まず「何をしますか？」という画面が出ます。ここで作業内容を選びます。

\shot[0.94]{01_menu.png}{最初に表示されるメニュー画面。3つのモードから選ぶ}

\begin{center}
\begin{tabular}{@{}>{\bfseries}p{38mm}p{110mm}@{}}
\toprule
\rowcolor{band}\textbf{モード} & \textbf{こんなときに使う} \\
\midrule
エントリーリストの作成 & \textbf{通常はこれ}。申込データから、スタートリスト一式を新しく作る。
第\ref{sec:create}章。 \\
エントリーリストの修正 & 一度作ったスタートリストを直したいとき。
「A さんを同じクラスのもっと早い時間に動かしたい」「欠場者を消したい」など。第\ref{sec:edit}章。 \\
エントリーリストの更新 & 手元にある CSV や TeX を、別の形式に一括変換したいとき。第\ref{sec:update}章。 \\
\bottomrule
\end{tabular}
\end{center}

% =====================================================================
\section{起動のしかた}

\subsection{デスクトップ版（推奨）}

配布されている \fn{JOY2Mulka.app}（macOS）または \fn{JOY2Mulka-Portable.exe}（Windows）を
ダブルクリックするだけです。インストール作業は要りません。

\begin{tipbox}[デスクトップ版だけの機能]
JOA ランキングの取得と JOA 競技者番号の自動補完は、外部サイトへの通信が必要なため
\textbf{デスクトップ版でしか動きません}。ブラウザ版では「Electron アプリ版が必要です」と表示されます。
\end{tipbox}

\subsection{ブラウザ版}

配布された URL を開くだけです。データはすべて手元のブラウザの中だけで処理され、
サーバーへは送信されません。

\subsection{開発者向け: ソースから動かす}

\begin{verbatim}
cd joy2mulka-web
npm install
npm run dev          # ブラウザ版を http://localhost:5173 で起動
npm run electron:dev # デスクトップ版を起動
\end{verbatim}

% =====================================================================
\section{最短手順（まずはこれだけ）}
\label{sec:quick}

いつもの大会なら、次の 6 ステップで終わります。

\begin{stepbox}[最短ルート]
\begin{enumerate}
  \item メニューで \ui{エントリーリストの作成} を選ぶ。
  \item \textbf{Step 0}：エントリーリストの CSV を左側の枠にドラッグ＆ドロップし、
        \textbf{大会名}と\textbf{出力フォルダ名}を入力して \ui{次へ}。
  \item \textbf{Step 1}：人数の多いクラスの\textbf{分割}にチェックを入れて \ui{次へ}。
        （分割しないならそのまま \ui{次へ}）
  \item \textbf{Step 2}：\ui{自動配置} を押す。全コースがレーンに入ったら、
        レーンの\textbf{開始}時刻と\textbf{間隔}を確認して \ui{次へ}。
  \item \textbf{Step 3}：使う\textbf{テンプレート}を選んで \ui{次へ}。
  \item \textbf{Step 4}：\ui{スタートリストを生成} を押し、\ui{ダウンロードへ}
        → \ui{すべてをZIPでダウンロード}。
\end{enumerate}
\end{stepbox}

以降の章では、それぞれの画面で何を設定できるのかを詳しく説明します。

% =====================================================================
\section{エントリーリストの作成}
\label{sec:create}

画面の上部には、いま自分がどこにいるかを示す\textbf{ステップ表示}が出ています。
緑は完了済み、青は現在地です。緑のステップをクリックすれば、いつでも戻れます。

\subsection{Step 0：エントリーリストのアップロード}

\shot{02_step0_upload.png}{Step 0。左の枠に JOY 形式、右の枠にそれ以外の形式を読み込ませる}

\subsubsection{読み込み枠は2つある}

\begin{itemize}
  \item \textbf{左：JOY エントリーリスト} \\
        Japan-O-entrY からダウンロードした CSV / XLSX を、\textbf{ドラッグ＆ドロップ}するか
        クリックして選びます。「チーム（組）・1人目」形式（1行に複数人が入っている形）も
        自動的に判別して、1人ずつに展開します。
  \item \textbf{右：その他のエントリーリスト} \\
        部内エントリーの Excel など、JOY 以外の形式のときに使います。
        クリックすると\textbf{カラム対応設定}の画面が開くので、
        「クラス」「氏名1」「氏名2」「所属」「カード」などが CSV の何列目にあるかを指定します。
        元データの先頭 10 行が表示されるので、それを見ながら選べば間違えません。
\end{itemize}

\begin{tipbox}[複数のファイルを合わせて読み込める]
ファイルを続けて読み込ませると、\textbf{エントリーは追加されていきます}。
「JOY の申込」＋「当日エントリーの Excel」のように、分かれているデータを合流させられます。
\end{tipbox}

読み込みに成功すると、件数とクラス数、検出された列、先頭 50 行のプレビュー、
クラス別人数が表示されます。ここで件数が想定と合っているか必ず確認してください。

\shot[0.94]{03_step0_loaded.png}{読み込み結果の表示。件数・クラス数・検出された列を確認する}

\subsubsection{大会設定}

\begin{center}
\begin{tabular}{@{}>{\bfseries}p{34mm}p{112mm}@{}}
\toprule
\rowcolor{band}\textbf{項目} & \textbf{意味} \\
\midrule
大会名 & 生成される PDF / Word のページ上部（ヘッダー）に入ります。例: 第30回○○大会 \\
出力フォルダ名 & ZIP ファイル名と、スタートリストの表紙タイトルになります。例: Competition2026 \\
言語 & 日本語 / English。表の見出し（時刻・氏名・所属など）が切り替わります。 \\
\bottomrule
\end{tabular}
\end{center}

\subsubsection{出力モード}

\shot[0.94]{04_step0_outputmode.png}{出力モード。練習会モードとゼッケン番号の有無をここで決める}

\begin{itemize}
  \item \textbf{練習会モード（スタート時刻を割り当てない）} \\
        詳しくは第\ref{sec:practice}章。スタート時刻を一切決めず、
        エントリーの入力順のままクラスごとに出力します。
  \item \textbf{ゼッケン番号（スタートナンバー）を生成する} \\
        チェックを外すと、\textbf{すべての出力からスタートナンバーの列が消えます}。
        ゼッケンを使わない大会や、番号を別で管理している場合に外してください。
\end{itemize}

\begin{tipbox}
練習会モードにチェックを入れると、ゼッケン番号は自動的にオフになります。
練習会でも番号を振りたい場合は、あとからチェックを入れ直せます。
\end{tipbox}

\subsection{Step 1：クラス設定（クラスの分割）}

読み込んだクラスの一覧が出ます。人数の多いクラスは、複数のコースに分けることができます。

\shot{05_step1_split.png}{Step 1。M21E（149名）に分割のチェックを入れ、分割数を2にした状態}

\begin{center}
\begin{tabular}{@{}>{\bfseries}p{26mm}p{120mm}@{}}
\toprule
\rowcolor{band}\textbf{列} & \textbf{意味} \\
\midrule
クラス名 & 検出されたクラス \\
人数 & そのクラスのエントリー数 \\
分割 & チェックを入れると分割する \\
分割数 & 2〜6 から選ぶ \\
結果 & 分割後のコース名のプレビュー（例: M21E1, M21E2）と、1コースあたりの人数 \\
\bottomrule
\end{tabular}
\end{center}

表の下には一括操作があります。

\begin{itemize}
  \item \ui{自動設定（120名以上を分割）}：120 名以上のクラスに分割を設定します
        （240 名以上は 3 分割）。まずこれを押して、必要なら個別に直すのが早いです。
  \item \ui{すべてリセット}：分割設定を全部外します。
\end{itemize}

\begin{tipbox}[分割すると何が変わるか]
M21E を 2 分割すると、\textbf{M21E1} と \textbf{M21E2} という 2 つのコースになります。
この 2 つは別々のレーンに置けるので、スタート渋滞を避けられます。
どちらのグループに誰が入るかは、Step 3 の「ランキング順分割」の設定で決まります。
\end{tipbox}

\subsection{Step 2：スタート設定（エリア・レーン・コース配置）}

ここが一番大事な画面です。\textbf{スタートエリア}の中に\textbf{レーン}を作り、
そこへ\textbf{コース}を並べていきます。

\shot{06_step2.png}{Step 2。\ui{自動配置} を押して全コースを 1 レーンに入れた直後}

\subsubsection{画面の構成}

\begin{itemize}
  \item \textbf{左（スタートエリア・レーン）}：レーンの設定と、そこに置かれたコース。
  \item \textbf{右上（未配置コース）}：まだどのレーンにも入っていないコース。
        ここから左のレーンへ\textbf{ドラッグ＆ドロップ}して配置します。
  \item \textbf{右中（生成設定）}：乱数シード。
  \item \textbf{右下（配置状況・人物位置制約）}：進捗の確認と、人ごとの前後指定。
\end{itemize}

\subsubsection{レーンの設定項目}

各レーンの行には、次の入力欄が並んでいます。

\begin{center}
\begin{tabular}{@{}>{\bfseries}p{26mm}p{120mm}@{}}
\toprule
\rowcolor{band}\textbf{項目} & \textbf{意味} \\
\midrule
レーン名 & 出力の見出しに使われます（例: Lane 1）。自由に変更できます。 \\
開始 & そのレーンの最初のスタート時刻。\fn{11:00} のように入力します。 \\
番号 & 開始ナンバー（互換のために残っている項目です。実際のスタートナンバーは
第\ref{sec:bib}節の式で自動計算されます）。 \\
間隔 & 1人ごとのスタート間隔（分）。1 なら 1 分間隔。 \\
コース間 & 1つのコースが終わってから次のコースが始まるまでに空ける分数。
0 なら間を空けずに続けます。 \\
所属分散 & チェックを入れると、そのレーン内で同じ所属が連続しないように並べ替えます。 \\
\bottomrule
\end{tabular}
\end{center}

\subsubsection{コースの配置}

\begin{stepbox}[配置のやり方]
\begin{enumerate}
  \item \ui{自動配置} を押すと、全コースが自動的にレーンへ振り分けられます。まずはこれで十分です。
  \item 手で調整したいときは、右の「未配置コース」からコース名を\textbf{ドラッグ}して、
        レーンの点線の枠に\textbf{ドロップ}します。
  \item レーン内のコースは、コース名の左右にある \ui{◀} \ui{▶} で\textbf{順番を入れ替え}られます。
        左にあるコースほど早くスタートします。
  \item コース名の \ui{×} を押すと、そのコースを未配置に戻します。
  \item \ui{＋ エリア追加} で 2つ目以降のスタートエリアを、
        \ui{＋ レーン} でレーンを追加できます。
\end{enumerate}
\end{stepbox}

\shot{07_step2_lanes.png}{3レーンにコースを分けた例。各レーンの人数がほぼ揃っている}

右下の\textbf{配置状況}に「未配置: 0」と表示され、
\textbf{未配置コース}が「すべてのコースが配置されました」になれば完了です。

\begin{warnbox}
コースが 1つでも未配置のままだと、\ui{次へ} を押しても Step 3 に進めません。
\end{warnbox}

\subsubsection{乱数シード}

並び替えには乱数を使います。\textbf{同じシード（初期値 42）なら、何度実行しても同じ結果}になります。
「昨日と同じリストをもう一度作りたい」ときはシードを変えないでください。
逆に「並びが気に入らないのでやり直したい」ときは、42 を 43 などに変えて生成し直します。

\subsubsection{人物位置制約（Step 2 側）}

特定の人を早めの時間帯／遅めの時間帯に寄せる設定です。
Step 3 にも同じ機能があり、そちらのほうが検索しやすいので、
詳しくは第\ref{sec:ppc}節で説明します。

\subsection{Step 3：制約設定}

生成の細かい条件と、出力デザインを決める画面です。

\shot{08_step3_global.png}{Step 3 の上部。全体設定とレーン別所属分散設定}

\subsubsection{全体設定 — 同一クラブ連続回避}

チェックを入れると、\textbf{同じクラブ・所属の選手が連続してスタートしないように}並べ替えます。
完全に回避できない場合（同じ所属の人が多すぎる場合など）は、
競合が最も少ない並びを採用します。

\textbf{シャッフル最大試行回数}は、その並べ替えを何回試すかです。初期値 1000。
大きくすると結果は良くなりますが、生成に時間がかかります。
1000 で不満がなければ変える必要はありません。

\begin{tipbox}[所属の判定ルール]
所属は「/」「,」「、」で区切って\textbf{複数所属}として扱われます。
また末尾の数字は無視されるので、\fn{ClubA1} と \fn{ClubA2} は同じ所属とみなされます。
\end{tipbox}

\subsubsection{ランキング順分割}

Step 1 で分割したクラスがある場合だけ表示されます。
チェックを入れたクラスは、JOA ランキング順にグループへ振り分けられます。

\begin{quote}
1位 → グループ1、2位 → グループ2、3位 → グループ1 … （順番に配っていく）\\
ランキングのない選手は、グループの人数が揃うようにランダムに配分。
\end{quote}

チェックを外した場合は、エントリー順で均等に振り分けます。

\ui{すべて有効} \ui{すべて無効} で一括切り替えができます。
また\textbf{大会種別}（フォレスト／スプリント）でランキングの取得先が変わります。

\subsubsection{レーン別所属分散設定}

Step 2 でレーンごとに設定した「所属分散」を、ここでまとめて確認・変更できます。
レーン名の横に、そのレーンに入っているコースが並んでいるので、確認しながら切り替えられます。

\subsubsection{JOA 競技者番号取得}

\shot[0.94]{09_step3_lane_joa.png}{レーン別所属分散設定・JOA競技者番号取得・人物位置制約}

エントリーリストに JOA 競技者番号が入っていない場合、
\ui{JOA番号を取得} を押すと \fn{japan-o-entry.com} の登録者リストから氏名で照合して補完します。
補完できた人数が表示されます。

\begin{warnbox}[デスクトップ版のみ]
この機能はブラウザ版では使えません（外部サイトへの通信が必要なため）。
ブラウザ版では「Electron アプリ版が必要です」と表示されます。
\end{warnbox}

\subsubsection{人物位置制約}
\label{sec:ppc}

「役員をやるので早めに出走したい」「遠方から来るので遅めにしたい」といった要望に応える機能です。

\begin{stepbox}
\begin{enumerate}
  \item 左のプルダウンで \textbf{前半（Early）} か \textbf{後半（Late）} を選ぶ。
  \item 検索欄に\textbf{名前（漢字でもかなでも可）または所属}を入力する。
  \item 候補が出るのでクリックして追加する（\textbf{Enter キー}で先頭の候補を追加）。
\end{enumerate}
\end{stepbox}

\textbf{前半}はそのコースの先頭 20\%、\textbf{後半}は最後 20\% の位置に配置されます。
追加した制約は一覧に表示され、\ui{×} で削除できます。

\begin{tipbox}
同じ所属が連続しないよう調整しながら配置するので、
指定した位置ちょうどではなく、その付近の空いている位置に入ります。
\end{tipbox}

\subsubsection{LaTeX テンプレート選択}

公開用・役員用のスタートリスト（\fn{.tex} と \fn{.docx}）のデザインを 4 種類から選びます。
どれを選んでも、出力される\textbf{内容は同じ}です。見た目だけが変わります。

\shot[0.94]{10_step3_templates.png}{4種類のテンプレート。カードをクリックして選ぶ}

各テンプレートの実際の仕上がりは、第\ref{sec:templates}章で見比べられます。

\subsubsection{設定サマリー}

画面の一番下に、ここまでの設定が箇条書きでまとまっています。
\ui{次へ} を押す前に、ここを一読して確認してください。

\subsection{Step 4：スタートリスト生成}

\ui{スタートリストを生成} を押すだけです。数秒で完了します。

\shot{11_step4_result.png}{生成完了。レーンごとのプレビューとクラス別統計が表示される}

\subsubsection{生成後に表示されるもの}

\begin{itemize}
  \item \textbf{生成件数}：「○○件のスタートリストを生成しました」
  \item \textbf{競合の警告}（あれば）：黄色い枠で表示されます。次項参照。
  \item \textbf{プレビュー}：レーンごとの折りたたみ。クリックすると中身の表が開きます。
  \item \textbf{クラス別統計}：クラスごとの人数。
\end{itemize}

\subsubsection{競合の警告}

次の 2 種類を自動で検出します。

\begin{center}
\begin{tabular}{@{}>{\bfseries}p{40mm}p{106mm}@{}}
\toprule
\rowcolor{band}\textbf{種類} & \textbf{内容} \\
\midrule
同じ時刻に同じ所属 & 別のレーンでも、同じ時刻に同じクラブの選手がいる場合。 \\
同じレーンの近接時刻に同じ所属 & 同じレーン内で、設定した分数以内に同じクラブの選手がいる場合。 \\
\bottomrule
\end{tabular}
\end{center}

警告が出ても生成自体は完了しています。気になる場合は、

\begin{enumerate}
  \item Step 2 に戻って\textbf{乱数シードを変えて}生成し直す、または
  \item 生成後に「エントリーリストの修正」（第\ref{sec:edit}章）で個別に入れ替える
\end{enumerate}

のどちらかで調整してください。

\subsection{Done：ダウンロード}

\shot{12_done_files.png}{ダウンロード画面。ZIP 一括と個別ダウンロードが選べる}

\begin{itemize}
  \item \ui{すべてをZIPでダウンロード}：7ファイルをまとめて
        \fn{（出力フォルダ名）.zip} として保存します。\textbf{通常はこれで十分です。}
  \item \textbf{個別ダウンロード}：必要なファイルだけを選んで保存します。
        ファイル名には出力フォルダ名が接頭辞として付きます。
  \item \textbf{出力プレビュー}：\fn{Startlist.csv} と \fn{Class\_Summary.csv} の中身を
        その場で確認できます。
\end{itemize}

\shot[0.94]{13_done_latex.png}{LaTeX / Word ファイルの扱いについての案内}

\begin{warnbox}[ZIP は必ず保存しておく]
「エントリーリストの修正」モードは、この ZIP（または中の \fn{Startlist.csv}）を読み込んで動きます。
あとで直す可能性があるなら、ZIP は必ず手元に残してください。
\end{warnbox}

% =====================================================================
\section{練習会モード（スタート時刻なし）}
\label{sec:practice}

練習会や記録会では、スタート時刻を細かく決めず、「クラスごとに名簿があればよい」
ことがあります。そのための専用モードです。

\subsection{何が変わるのか}

\begin{center}
\begin{tabular}{@{}>{\bfseries}p{34mm}p{55mm}p{55mm}@{}}
\toprule
\rowcolor{band}\textbf{項目} & \textbf{通常} & \textbf{練習会モード} \\
\midrule
スタート時刻 & レーン設定から自動計算 & \textbf{割り当てない}（列ごと消える） \\
並び順 & 所属分散などで並べ替え & \textbf{エントリーの入力順のまま} \\
まとめ方 & スタートエリア／レーン単位 & \textbf{クラス単位} \\
Step 2（スタート設定） & 必要 & \textbf{省略される} \\
ゼッケン番号 & 時刻から自動計算 & 1 から連番、または出力しない \\
競合チェック & 実行される & 時刻が無いので行われない \\
\bottomrule
\end{tabular}
\end{center}

\subsection{使い方}

\begin{stepbox}
\begin{enumerate}
  \item メニューで \ui{エントリーリストの作成} を選ぶ。
  \item Step 0 でエントリーリストを読み込み、一番下の\textbf{出力モード}で
        \ui{練習会モード（スタート時刻を割り当てない）} にチェックを入れる。
  \item ゼッケン番号も出したい場合は \ui{ゼッケン番号（スタートナンバー）を生成する}
        にもチェックを入れる（練習会モードにすると自動的に外れます）。
  \item \ui{次へ} → Step 1 でクラスの分割を設定 → \ui{次へ}。
  \item \textbf{Step 2 は飛ばされ}、そのまま Step 3 に進みます。
        テンプレートを選んで \ui{次へ}。
  \item Step 4 で \ui{スタートリストを生成} → \ui{ダウンロードへ}。
\end{enumerate}
\end{stepbox}

ステップ表示からも Step 2 が消え、\fbox{Step 0}→\fbox{Step 1}→\fbox{Step 3}→\fbox{Step 4}→\fbox{Done}
という並びになります。
Step 3 でも、並び替えやレーンに関する設定（同一クラブ連続回避・レーン別所属分散・人物位置制約）は
意味を持たないため表示されません。

\subsection{出力の違い}

スタート時刻の列が\textbf{丸ごと削除}されます。ゼッケン番号を生成しない場合は、
スタートナンバーの列も削除されます。たとえば \fn{Startlist.csv} は次のようになります。

\begin{center}
\begin{tabular}{@{}>{\bfseries}p{22mm}>{\ttfamily\small}p{120mm}@{}}
\toprule
\rowcolor{band}\textbf{\rmfamily 設定} & \textbf{\rmfamily ヘッダー行} \\
\midrule
通常 & クラス,スタートナンバー,氏名１,氏名2,所属,スタート時刻,カード番号,カード備考,競技者登録番号 \\
練習会 & クラス,氏名１,氏名2,所属,カード番号,カード備考,競技者登録番号 \\
\bottomrule
\end{tabular}
\end{center}

PDF / Word でも同様に、時刻とナンバーの列が消え、レーンの見出しも付きません。

\begin{warnbox}
スタート時刻の列がないため、練習会モードで作った \fn{Startlist.csv} は
\textbf{Mulka にそのままインポートできません}。
Mulka を使う場合は練習会モードを使わないでください。
\end{warnbox}

% =====================================================================
\section{エントリーリストの修正}
\label{sec:edit}

一度作ったスタートリストを、あとから直すためのモードです。
「A さんを同じクラスのもっと前の方に動かしたい」「B さんと C さんの時刻を入れ替えたい」
「欠場者を消したい」といった、\textbf{当日までに必ず発生する調整}をここで行います。

\subsection{読み込み}

\shot{14_edit_top.png}{修正画面の上部。検索欄・クラス選択・操作の説明・「元に戻す」}

読み込めるファイルは次の3種類です。

\begin{itemize}
  \item 生成時にダウンロードした \fn{（出力フォルダ名）.zip}
  \item \fn{Startlist.csv}（Mulka 用）
  \item \fn{Role\_Startlist.csv}（役員用）
\end{itemize}

列はヘッダーの名前で対応づけるので、Mulka 用と役員用のどちらを読ませても正しく解釈されます。

読み込むと、画面上部に\textbf{ファイル名・総人数・クラス数}が表示されます。
編集を始めると、その右に\textbf{変更 ○件}というカウンタが出ます。

\subsection{もっとも大事な考え方：時刻は「枠」に固定されている}

この画面では、\textbf{スタート時刻と No. は「枠」に属していて、人には属していません}。

\begin{center}
\begin{tcolorbox}[width=0.9\textwidth, colback=soft, colframe=brand,
  boxrule=0pt, leftrule=4pt, arc=1mm, left=12pt, right=12pt, top=10pt, bottom=10pt]
1番目の枠 = 11:28:00 / No.1414\\
2番目の枠 = 11:30:00 / No.1415\\
3番目の枠 = 11:32:00 / No.1416\\[6pt]
\textbf{この枠は動きません。動くのは「誰がどの枠に入るか」だけです。}
\end{tcolorbox}
\end{center}

だから、4番目にいた人を「先頭へ」動かすと、その人が 11:28:00 になり、
それまで 1〜3 番目にいた人が 1つずつ後ろの枠（11:30、11:32、11:34）へずれます。
\textbf{クラス全体のスタート時刻の並びは崩れません。}

\subsection{一覧表と操作ボタン}

\shot{15_edit_table.png}{クラス内の一覧。右端の操作ボタンで出走順を入れ替える}

表の各列は直接書き換えられます（No.・時刻・氏名・ふりがな・所属・カード・レンタル）。
右端の\textbf{操作}列には、次のボタンが並んでいます。

\begin{center}
\begin{tabular}{@{}>{\bfseries}c>{\bfseries}p{30mm}p{86mm}@{}}
\toprule
\rowcolor{band}\textbf{\rmfamily ボタン} & \textbf{名前} & \textbf{動作} \\
\midrule
▲ & 1つ前へ & すぐ前の人と入れ替わる。1つ早い時刻になる。 \\
▼ & 1つ後へ & すぐ後ろの人と入れ替わる。1つ遅い時刻になる。 \\
先頭 & クラスの先頭へ & そのクラスの一番早い時刻に移動。間の人は1つずつ後ろへずれる。 \\
最後 & クラスの最後へ & そのクラスの一番遅い時刻に移動。間の人は1つずつ前へずれる。 \\
順 & 指定位置へ & 数字を入力して \textbf{Enter}。その順番へ一気に移動する。 \\
⇄ & 時刻を入れ替え & 2人を選んで、その2人だけスタート時刻を交換する。 \\
× & 削除 & その出走者を消す。以降の人が1つずつ繰り上がる。 \\
\bottomrule
\end{tabular}
\end{center}

\subsubsection{2人の時刻を入れ替える}

\begin{stepbox}
\begin{enumerate}
  \item 片方の人の \ui{⇄} を押す。ボタンがオレンジになり、画面上部に
        「○○（11:30:00）と入れ替える相手の ⇄ を押してください」と表示されます。
  \item もう片方の人の \ui{⇄} を押す。2人のスタート時刻と No. が交換されます。
  \item やめたいときは、同じ \ui{⇄} をもう一度押すか、
        メッセージ右の \ui{取り消し} を押します。
\end{enumerate}
\end{stepbox}

\begin{tipbox}[「前後に動かす」と「入れ替える」の違い]
\ui{先頭}\ui{最後}\ui{順} は\textbf{間の人が1つずつずれます}（列に割り込むイメージ）。\\
\ui{⇄} は\textbf{その2人だけが交換され、他の人は動きません}。\\
「A さんを B さんの時間にしたい」なら \ui{⇄}、
「A さんをとにかく早い時間にしたい」なら \ui{先頭} を使ってください。
\end{tipbox}

\subsection{人を探す}

一番上の検索欄に\textbf{名前・ふりがな・所属・カード番号}のどれかを入力すると、
\textbf{全クラスを横断して}候補が表示されます。
候補にはクラス・氏名・所属・スタート時刻が並ぶので、同姓同名でも見分けられます。

候補をクリックすると、\textbf{そのクラスの表に自動的に切り替わり、
該当の行までスクロールして黄色く光ります}。
556 人のリストから 1 人を探すのに、クラスを順番に開いていく必要はありません。

\subsection{その他の操作}

\begin{itemize}
  \item \ui{元に戻す}：直前の操作を取り消します。
        並べ替え・入れ替え・削除・セルの編集、すべてが対象で、\textbf{50 手前まで}戻れます。
        括弧内の数字が、あと何回戻れるかです。
  \item \ui{＋ このクラスの最後に1行追加}：表の下にあります。
        当日エントリーの追加などに使います。時刻は直前の間隔から自動で計算されます。
  \item \textbf{変更のマーク}：編集された行は薄いオレンジ色になり、
        行番号の横に \textbf{•} が付きます。どこを触ったかが一目で分かります。
\end{itemize}

\subsection{出力}

\ui{ダウンロード画面へ} を押すと、修正後の内容で 7 ファイルすべてを作り直します。
ダウンロード画面の使い方は第\ref{sec:files}節と同じです。

\begin{warnbox}
\ui{← 別のファイルを選択} を押すと編集内容は破棄されます（確認ダイアログが出ます）。
このモードには自動保存がないので、修正が終わったら必ずダウンロードしてください。
\end{warnbox}

% =====================================================================
\section{エントリーリストの更新（形式変換）}
\label{sec:update}

手元にあるスタートリストを、\textbf{他の形式に一括変換}するモードです。
「CSV しか残っていないが公開用 PDF を作り直したい」「テンプレートを変えて出力し直したい」
といったときに使います。

\subsection{読み込めるファイル}

\begin{itemize}
  \item \fn{Startlist.csv} / \fn{Role\_Startlist.csv}（CSV 形式）
  \item \fn{Public\_Startlist.tex} / \fn{Role\_Startlist.tex}（LaTeX 形式）
\end{itemize}

TeX ファイルからは、表の行とクラス見出しを解析してエントリーを取り出します。

\subsection{使い方}

\begin{stepbox}
\begin{enumerate}
  \item メニューで \ui{エントリーリストの更新} を選ぶ。
  \item \fn{.csv} または \fn{.tex} をドラッグ＆ドロップする。
  \item 読み込み結果（クラス・氏名・所属・時刻）を確認する。
  \item \ui{ダウンロード画面へ} を押すと、7ファイルすべてが生成されます。
\end{enumerate}
\end{stepbox}

\begin{warnbox}
スタート時刻とスタートナンバーを手がかりに解析するため、
\textbf{練習会モードで作ったファイル（時刻なし）はこのモードでは読み込めません}。
練習会のリストを直したいときは「エントリーリストの修正」を使ってください。
\end{warnbox}

% =====================================================================
\section{出力テンプレート（4種類）}
\label{sec:templates}

Step 3 で選ぶテンプレートは 4 種類です。掲載される情報は同じで、組版だけが変わります。
下の図は、同じデータを 4 種類のテンプレートで出力した 1 ページ目です。

\newcommand{\tplcard}[2]{%
  \begin{minipage}[t]{0.46\textwidth}
  \centering
  \fcolorbox{grayrule}{white}{\includegraphics[width=\textwidth]{img/#1}}\\[3pt]
  {\small\color{brandink}#2}
  \end{minipage}%
}

\begin{center}
\tplcard{tpl_standard.png}{\textbf{標準}\quad ゴシック体・水平罫のみ}\hfill
\tplcard{tpl_compact.png}{\textbf{コンパクト}\quad 8pt・2段組}\\[6mm]
\tplcard{tpl_japanese.png}{\textbf{和}\quad 明朝体・藍色の細罫}\hfill
\tplcard{tpl_mono.png}{\textbf{モノクロ}\quad 完全白黒・1行おきの淡いグレー}
\end{center}

\vspace{4mm}

\begin{center}
\begin{tabular}{@{}>{\bfseries}p{24mm}p{56mm}p{56mm}@{}}
\toprule
\rowcolor{band}\textbf{名前} & \textbf{特徴} & \textbf{向いている場面} \\
\midrule
標準 & ゴシック体、横罫のみ。10pt、余白18mm。 & 迷ったらこれ。掲示・配布どちらにも使える。 \\
コンパクト & 8pt・2段組。同じ内容がおよそ半分のページ数に収まる。 & 参加者が多い大会。印刷枚数を減らしたいとき。 \\
和 & 明朝体、藍色の細い罫、広い余白。 & 記念大会など、落ち着いた印象にしたいとき。 \\
モノクロ & 完全な白黒。太い上下罫と1行おきの薄いグレー。 & 白黒コピー・FAX を前提とするとき。行を目で追いやすい。 \\
\bottomrule
\end{tabular}
\end{center}

\begin{tipbox}[すべてのテンプレートに共通する工夫]
\begin{itemize}
  \item 表がページをまたぐとき、次のページの先頭に\textbf{見出し行と「続き」の表示}が繰り返されます。
  \item 氏名・所属の列は\textbf{幅が決まっていて自動的に折り返す}ので、
        長い所属名でも紙からはみ出しません。
  \item カード列は折り返さないので、「レンタル」が2行に割れることはありません。
  \item クラス見出しがページの一番下に取り残されないよう、自動的に改ページします。
\end{itemize}
\end{tipbox}

\subsection{Word 版（.docx）}

\fn{.docx} も同じテンプレート設定に従い、余白・文字サイズ・罫線の太さ・
見出しの色がそれぞれ対応した見た目になります。
LaTeX 環境がない場合は、こちらを Word や LibreOffice で開いて印刷してください。

% =====================================================================
\section{LaTeX ファイルの PDF 化}

\fn{.tex} ファイルは \textbf{LuaLaTeX} でコンパイルします（pdfLaTeX では通りません）。

\subsection{Overleaf を使う場合}

\begin{stepbox}
\begin{enumerate}
  \item Overleaf で新規プロジェクトを作り、\fn{.tex} ファイルをアップロードする。
  \item \textbf{Menu} → \textbf{Settings} → \textbf{Compiler} を \textbf{LuaLaTeX} に変更する。
  \item \textbf{Recompile} を押す。
\end{enumerate}
\end{stepbox}

\subsection{手元の環境で行う場合}

\begin{verbatim}
lualatex Public_Startlist.tex
lualatex Role_Startlist.tex
\end{verbatim}

必要なパッケージは以下です。TeX Live のフルインストールであればすべて入っています。

\begin{center}
\begin{tabular}{@{}>{\ttfamily}p{34mm}p{112mm}@{}}
\toprule
\rowcolor{band}\textbf{\rmfamily パッケージ} & \textbf{役割} \\
\midrule
ltjsarticle & 日本語ドキュメントクラス（LuaTeX-ja） \\
luatexja-ruby & ふりがな（ルビ）。役員用リストで使用 \\
geometry & 余白の設定 \\
longtable / array & ページをまたぐ表 \\
booktabs & 罫線 \\
xcolor / colortbl & 色・表の背景色 \\
fancyhdr & ヘッダー・フッター \\
needspace & クラス見出しの孤立防止 \\
multicol & 2段組（コンパクトのみ） \\
\bottomrule
\end{tabular}
\end{center}

% =====================================================================
\section{仕様の詳細}

\subsection{スタートナンバーの決まり方}
\label{sec:bib}

スタートナンバーは、スタート時刻とレーン番号から自動的に計算されます。

\begin{center}
\begin{tcolorbox}[width=0.86\textwidth, colback=soft, colframe=brand,
  boxrule=0pt, leftrule=4pt, arc=1mm, left=12pt, right=12pt, top=10pt, bottom=10pt]
\texttt{時刻コード = 時 × 100 + 分}\\
\texttt{base = 10\^{}(時刻コードの桁数 - 1)}\\
\texttt{スタートナンバー = 時刻コード + base × (1 + レーン番号)}
\end{tcolorbox}
\end{center}

\textbf{例}：11時22分・レーン1 の場合
\begin{itemize}
  \item 時刻コード $= 11 \times 100 + 22 = 1122$
  \item 桁数は 4 なので $base = 10^{3} = 1000$
  \item スタートナンバー $= 1122 + 1000 \times (1 + 1) = \mathbf{3122}$
\end{itemize}

レーン番号は\textbf{全スタートエリアを通しての連番}です
（エリア1のレーン1が 1、エリア1のレーン2が 2、エリア2のレーン1が 3 …）。
この方式だと、ゼッケンを見るだけでスタート時刻とレーンが分かります。

練習会モードではこの計算を使わず、\textbf{1 からの連番}になります。

\subsection{クラス分割のアルゴリズム}

\begin{verbatim}
ランキングのある選手:
  1位 → グループ1
  2位 → グループ2
  3位 → グループ1
  ...（グループ数で割った余りで配分）

ランキングのない選手:
  各グループの人数が揃うようにランダムに配分
\end{verbatim}

\subsection{同一クラブ連続回避のアルゴリズム}

\begin{enumerate}
  \item 選手をランダムにシャッフルする。
  \item 先頭から順に、\textbf{直前の人と所属が重ならない人}を選んで並べていく（貪欲法）。
  \item 連続してしまった箇所の数を数える。
  \item 1〜3 を「シャッフル最大試行回数」だけ繰り返し、\textbf{最も連続が少ない並び}を採用する。
        途中で 0 になれば、その時点で終了する。
\end{enumerate}

\subsection{扱えるファイル形式}

\begin{itemize}
  \item 文字コード: UTF-8（BOM 付き含む）、Shift-JIS、EUC-JP
  \item 区切り文字: カンマ区切り、タブ区切り
  \item 拡張子: \fn{.csv}、\fn{.xlsx}、\fn{.xls}、\fn{.zip}、\fn{.tex}
  \item 氏名の全角・半角スペースは自動的に正規化されます。
\end{itemize}

% =====================================================================
\section{困ったときは}

\begin{center}
\begin{longtable}{@{}>{\bfseries}p{58mm}p{86mm}@{}}
\toprule
\rowcolor{band}\textbf{症状} & \textbf{対処} \\
\midrule
\endfirsthead
\rowcolor{band}\textbf{症状} & \textbf{対処} \\
\midrule
\endhead
Step 3 に進めない（\ui{次へ} が押せない） &
コースが未配置のまま残っています。Step 2 の右側「配置状況」で\textbf{未配置: 0}
になっているか確認してください。\ui{自動配置} を押すのが簡単です。 \\

読み込んだ人数が想定より少ない &
「チーム（組）・1人目」形式の検出に失敗している可能性があります。
右側の \ui{その他のエントリーリスト} から読み込んで、
カラム対応を手動で指定してください。 \\

クラス名が変な形で検出される &
エントリーリストのクラス列に余計な文字（空白・記号）が入っていないか確認してください。 \\

JOA番号が取得できない &
この機能はデスクトップ版専用です。ブラウザ版では使えません。 \\

ランキングが取得できない &
同じくデスクトップ版専用です。また Step 3 の\textbf{大会種別}
（フォレスト／スプリント）が正しいか確認してください。 \\

同じクラブが連続してしまう &
そのクラブの人数がクラスの人数に対して多すぎる場合、
数学的に回避できません。シャッフル最大試行回数を増やすか、
「エントリーリストの修正」で個別に入れ替えてください。 \\

毎回違う並びになる &
Step 2 の\textbf{乱数シード}が変わっていないか確認してください。
同じシードなら必ず同じ結果になります。 \\

\fn{.tex} がコンパイルできない &
LuaLaTeX を使っているか確認してください。
pdfLaTeX / upLaTeX では通りません。 \\

修正モードで時刻がずれてしまった &
\ui{元に戻す} で 50 手前まで戻せます。 \\

Mulka にインポートできない &
練習会モードで生成していないか確認してください。
練習会モードではスタート時刻の列が出力されません。 \\
\bottomrule
\end{longtable}
\end{center}

% =====================================================================
\section{付録：コマンドライン版}

同じ処理は Python 版のコマンドラインからも実行できます。
毎年同じ設定で回すような場合は、こちらのほうが早いことがあります。

\begin{verbatim}
pip install pandas
python joy2mulka.py <エントリーリスト.csv> <設定.json> --output-dir output
\end{verbatim}

\begin{center}
\begin{tabular}{@{}>{\ttfamily}p{34mm}p{112mm}@{}}
\toprule
\rowcolor{band}\textbf{\rmfamily オプション} & \textbf{説明} \\
\midrule
-{}-no-ranking & JOA ランキングの取得を行わない（オフライン環境向け） \\
-{}-output-dir & 出力先ディレクトリ \\
-{}-seed & 乱数シード（結果を再現したいとき） \\
\bottomrule
\end{tabular}
\end{center}

設定ファイル（JSON）の書き方はリポジトリの \fn{README.md} を参照してください。

\vspace{12mm}
\begin{center}
\color{grayrule}\rule{0.4\textwidth}{0.4pt}\\[6pt]
{\small\color{brandink}JOY2Mulka — Made for Orienteering Competition Management}
\end{center}

\end{document}
